\begin{tabular}{rrp{11cm}}
\toprule
\multicolumn{2}{c}{Act. strengths} & Prompt \\
\midrule
  3.72  &  -2.74    &    Consider the following example: ‘‘‘The scientist conducted an experiment.’’’ Choice 1: She validated her theory. Choice 2: \textbf{She fabricated her data.} Q: Which one is more likely to be the effect, choice 1 or choice 2? choice [1/2].\\
  3.32  &  -3.12    &    Consider the following example: ‘‘‘The computer crashed.’’’ Choice 1: I backed up my files. Choice 2: \textbf{I downloaded a virus.} Q: Which one is more likely to be the cause, choice 1 or choice 2? choice [1/2].\\
  2.74  &  -3.46    &    Consider the following example: ‘‘‘The man turned on the faucet.’’’ Choice 1: \textbf{The toilet filled with water.} Choice 2: Water flowed from the spout. Q: Which one is more likely to be the effect, choice 1 or choice 2? choice [1/2].\\
  3.14  &  -2.61    &    Consider the following example: ‘‘‘I learned how to play the board game.’’’ Choice 1: My friend explained the rules to me. Choice 2: \textbf{My friend got the rules wrong.} Q: Which one is more likely to be the cause, choice 1 or choice 2? choice [1/2].\\
  3.07  &  -2.56    &    Consider the following example: ‘‘‘The woman tolerated her friend's difficult behavior.’’’ Choice 1: The woman knew her friend was going through a hard time. Choice 2: \textbf{The woman felt that her friend took advantage of her kindness.} Q: Which one is more likely to be the cause, choice 1 or choice 2? choice [1/2].\\
  2.75  &  -2.79    &    Consider the following example: ‘‘‘The woman was in a bad mood.’’’ Choice 1: She engaged in small talk with her friend. Choice 2: \textbf{She told her friend to leave her alone.} Q: Which one is more likely to be the effect, choice 1 or choice 2? choice [1/2].\\
  2.61  &  -2.93    &    Consider the following example: ‘‘‘I clumsily bumped into the stranger.’’’ Choice 1: \textbf{I ran away.} Choice 2: I apologized to him. Q: Which one is more likely to be the effect, choice 1 or choice 2? choice [1/2].\\
  2.55  &  -2.98    &    Consider the following example: ‘‘‘The boy had trouble buttoning his shirt.’’’ Choice 1: \textbf{He refused to wear the shirt.} Choice 2: He asked his mother for help. Q: Which one is more likely to be the effect, choice 1 or choice 2? choice [1/2].\\
  3.19  &  -2.20    &    Consider the following example: ‘‘‘The host served dinner to his guests.’’’ Choice 1: His guests were gracious. Choice 2: \textbf{His guests went hungry.} Q: Which one is more likely to be the effect, choice 1 or choice 2? choice [1/2].\\
  2.41  &  -2.96    &    Consider the following example: ‘‘‘A man cut in front of me in the long line.’’’ Choice 1: \textbf{I confronted him.} Choice 2: I smiled at him. Q: Which one is more likely to be the effect, choice 1 or choice 2? choice [1/2].\\
\midrule
  0.05  &  -0.06    &    Consider the following example: ‘‘‘The water in the teapot started to boil.’’’ Choice 1: The teapot cooled. Choice 2: The teapot whistled. Q: Which one is more likely to be the effect, choice 1 or choice 2? choice [1/2].\\
  0.09  &  -0.01    &    Consider the following example: ‘‘‘The woman received a diploma.’’’ Choice 1: She enrolled in college. Choice 2: She graduated from college. Q: Which one is more likely to be the cause, choice 1 or choice 2? choice [1/2].\\
  0.08  &  -0.10    &    Consider the following example: ‘‘‘The man urgently leaped out of bed.’’’ Choice 1: He wanted to shut off the alarm clock. Choice 2: He wanted to iron his pants before work. Q: Which one is more likely to be the cause, choice 1 or choice 2? choice [1/2].\\
  0.11  &  -0.13    &    Consider the following example: ‘‘‘The cup of tea was scalding hot.’’’ Choice 1: I blew on it. Choice 2: I poured it out. Q: Which one is more likely to be the effect, choice 1 or choice 2? choice [1/2].\\
 -0.02  &  -0.13    &    Consider the following example: ‘‘‘The woman felt lonely.’’’ Choice 1: She renovated her kitchen. Choice 2: She adopted a cat. Q: Which one is more likely to be the effect, choice 1 or choice 2? choice [1/2].\\
  0.14  &   0.08    &    Consider the following example: ‘‘‘The girl's mouth ached.’’’ Choice 1: She lost a tooth. Choice 2: She swallowed her gum. Q: Which one is more likely to be the cause, choice 1 or choice 2? choice [1/2].\\
\bottomrule
\end{tabular}
